% Fuentes exactas: eq:Jbeta y eq:Fbeta. Se inserta antes de las contracciones.
\begin{frame}[label=nav-frame-108]{Momentos de $Q_\beta$: gradiente y campo escalar}
\TablaSetup
\begin{tabular}{@{}p{.19\textwidth}p{.48\textwidth}p{.26\textwidth}@{}}
\TablaCabecera{Momento} & \TablaCabecera{Resultado obtenido} & \TablaCabecera{Dependencia} \\
\TablaLinea
$\Pi_\beta^a$ & $\ell^2\beta_0\big(6R^a{}_b\nabla^b\phi-2R\nabla^a\phi\big)$ & Respecto del gradiente $\nabla_a\phi$. \\[5mm]
$F_\phi^{(\beta)}$ & $0$ & Sin dependencia explícita en $\phi$. \\
\TablaLinea
\end{tabular}\par
\vspace{5mm}

Ya contamos con los cuatro momentos de $Q_\beta$.
El siguiente paso combina estos resultados para obtener
$\mathcal R^{(\beta)}_{(ab)}$, las ecuaciones y la frontera.
\end{frame}
